% mazzi: 09
\chapter{L'allenamento della potenza muscolare}

\textbf{Allenamento della potenza muscolare}: prescrizione ricavata dalle curve forza-velocità e
potenza-velocità della dinamometria isoinerziale, dove il riferimento non è più il numero di
ripetizioni ma il \textbf{carico} e la \textbf{potenza espressa}.

% ---------------------------------------------------------------------------
\section{I parametri dell'allenamento neuromuscolare}

Due soli parametri definiscono lo stimolo:
\begin{enumerate}
  \item \textbf{entità del carico}: in percentuale dell'1RM;
  \item \textbf{intensità dello stimolo}: in percentuale della potenza massima ($P_{max}$).
\end{enumerate}

\begin{itemize}
  \item la \textbf{potenza} è la vera intensità dell'allenamento: il numero di ripetizioni smette di
    essere il parametro di riferimento;
  \item una volta determinata la curva forza-velocità, sono noti i valori di potenza di ogni carico
    sotteso alla curva;
  \item il carico si sceglie rispetto al \textbf{picco di potenza} --- a sinistra o a destra ---
    secondo l'esigenza e l'obiettivo.
\end{itemize}

% ---------------------------------------------------------------------------
\section{I parametri di riferimento dei quattro tipi di allenamento}

Ogni tipo di allenamento occupa una porzione diversa delle due curve, individuata da una coppia di
valori: quanto carico (\% 1RM) e quanta potenza (\% $P_{max}$) si deve esprimere con quel carico.

\begin{itemize}
  \item \textbf{forza massima} (\textit{maximal strength training}): carico 70--100\,\% 1RM, potenza
    $>$ 90\,\% $P_{max}$;
  \item \textbf{ipertrofia} (\textit{hypertrophic training}): carico 70--90\,\% 1RM, potenza
    70--80\,\% $P_{max}$;
  \item \textbf{forza esplosiva} (\textit{explosive strength training}): carico 20--70\,\% 1RM,
    potenza $>$ 90\,\% $P_{max}$;
  \item \textbf{resistenza alla forza veloce} (\textit{speed strength endurance training}): carico
    20--50\,\% 1RM, potenza 80--90\,\% $P_{max}$.
\end{itemize}

\rubrica{Perché la potenza cambia il senso del carico}
\begin{itemize}
  \item con carichi $>$ 70\,\% 1RM il reclutamento è quasi massimale: a fine serie tutte le fibre
    risultano esaurite, ma quelle davvero reclutate e portate a esaurimento sono le \textbf{fibre
    veloci};
  \item nell'ipertrofia la potenza si abbassa deliberatamente (fino all'80\,\%) per stressare il
    sistema dal punto di vista metabolico: maggiore produzione di acido lattico $\rightarrow$
    aumento del GH $\rightarrow$ ipertrofia; interessa soprattutto lo stimolo sul
    \textbf{testosterone}, fenotipizzatore delle fibre veloci;
  \item nella forza esplosiva il carico può stare a destra come a sinistra del picco, ma la potenza
    richiesta resta $>$ 90\,\%: l'obiettivo è condizionare \textbf{pattern neuromuscolari}
    specifici, cioè strategie di reclutamento specifiche;
  \item nella resistenza alla forza veloce il carico si sbilancia verso destra (carichi bassi,
    velocità alte) per ottenere uno \textbf{stress metabolico}.
\end{itemize}

Il quadro di riferimento esteso aggiunge la \textbf{resistenza muscolare} e propone per la
resistenza alla forza veloce un intervallo di carico leggermente diverso:

\begin{table}[htbp]
  \centering
  \small
  \renewcommand{\arraystretch}{1.3}
  \begin{tabularx}{\textwidth}{|X|c|c|}
    \hline
    \textbf{Tipo di allenamento} & \textbf{Carico (\% 1RM)} &
      \textbf{\% della $P_{max}$ con quel carico} \\
    \hline
    Forza massima (\textit{maximal strength}) & 70--100\,\% & minimo 90\,\% \\
    \hline
    Forza esplosiva (\textit{speed/strength, power}) & 20--70\,\% & minimo 90\,\% \\
    \hline
    Ipertrofia (\textit{hypertrophy}) & 70--90\,\% & 75--85\,\% \\
    \hline
    Resistenza alla forza veloce (\textit{speed/endurance}) & 30--50\,\% & 80--90\,\% \\
    \hline
    Resistenza muscolare (\textit{muscle endurance}) & 30--70\,\% & 70--85\,\% \\
    \hline
  \end{tabularx}
\end{table}

% ---------------------------------------------------------------------------
\section{Reclutamento delle fibre e recupero tra le serie}

Le tre popolazioni di fibre --- \textbf{ST}, \textbf{FTa}, \textbf{FTb} --- si seguono durante e
dopo lo sforzo distinguendo le fibre \textit{recruited}, \textit{recruited exhausted}, \textit{not
recruited}, \textit{recovered} e \textit{not recovered}.

\rubrica{Durante la serie con carichi $>$ 70\,\% 1RM}
\begin{enumerate}
  \item \textbf{inizio}: sono reclutate le ST e parte delle FTa, le FTb solo in quota minima;
  \item \textbf{durante}: la quota di fibre reclutate ed esaurite cresce e si estende alle FTa e
    alle FTb;
  \item \textbf{fine}: quasi tutte le fibre risultano reclutate ed esaurite.
\end{enumerate}

Lo stesso avviene nella successione di \textbf{sforzi balistici massimali}: anche con carichi bassi
il reclutamento parte dalle fibre veloci e, man mano che si esauriscono, passa alle intermedie e poi
alle lente $\rightarrow$ calo dei valori di potenza.

\rubrica{Il recupero dopo ripetizioni submassimali}
\begin{itemize}
  \item \textbf{subito dopo il lavoro}: la potenza massima esprimibile è nettamente ridotta;
  \item \textbf{dopo 1 minuto}: le ST sono recuperate, ma le FTa risultano ancora \textit{not
    recovered} $\rightarrow$ la potenza non torna al 100\,\%;
  \item \textbf{dopo 3 minuti}: il recupero è più ampio, ma la quota di FTb non recuperate resta il
    fattore che limita la ripresa della serie.
\end{itemize}

Ne segue che il \textbf{tempo di recupero tra le serie} va commisurato alle fibre che interessano:
se sono le veloci, i tempi diventano lunghi e vanno verificati sul campo --- se dopo tre minuti la
potenza esprimibile resta $<$ 90\,\% della massima, il recupero si allunga ancora.

\rubrica{Ipertrofia selettiva}
\textbf{Ipertrofia selettiva}: condizionare solo le fibre che servono alla prestazione, attraverso i
pattern neuromuscolari altamente specifici legati alle \textbf{unità motorie veloci di tipo B}.
\begin{itemize}
  \item si contrappone all'\textbf{ipertrofia generale} del \textit{bodybuilder}, dove l'obiettivo è
    aumentare la massa a prescindere dal pattern neuromuscolare utilizzato;
  \item condizionare unità motorie che non sono il fattore limitante è inutile $\rightarrow$ può
    addirittura disturbare la prestazione finale.
\end{itemize}

I mezzi di allenamento si dispongono lungo un continuum tra \textbf{fenomeni nervosi} e
\textbf{massa muscolare}: i carichi leggeri ad alta velocità e le serie da 1--3RM stanno sul
versante nervoso, il 6RM in posizione intermedia e il 10RM verso la massa muscolare.

% ---------------------------------------------------------------------------
\section{Il monitoraggio della potenza nella serie}

Con l'encoder la potenza media di ogni ripetizione è disponibile in tempo reale. Nell'esempio di una
serie da 10 ripetizioni, la potenza media (W) segue l'andamento:

\begin{center}
246 --- 233 --- 245 --- 229 --- 219 --- 212 --- 196 --- 186 --- 175 --- 188
\end{center}

\begin{itemize}
  \item si fissa una \textbf{soglia di potenza} (per la forza esplosiva $\geq$ 90\,\% della
    $P_{max}$, con carichi indicativamente tra il 30--40\,\% e il 70\,\%);
  \item finché ogni ripetizione resta sopra la soglia si sta lavorando sulle unità motorie veloci;
  \item quando la potenza scende sotto il target è cambiato il pattern di reclutamento, perché
    quello privilegiato si è affaticato $\rightarrow$ la serie si interrompe;
  \item il \textbf{numero ottimale di ripetizioni} non è quindi il massimo eseguibile con quel
    carico, ma il numero di ripetizioni che rientrano nel 90\,\% della potenza massima;
  \item lo stesso criterio individualizza ripetizioni per serie, recupero tra le serie e numero di
    serie: il ``\textit{no pain no gain}'' non vale più, lo sforzo non deve essere massimo ma
    selettivo.
\end{itemize}

Poiché l'adattamento fa crescere le potenze espresse, i valori vanno rimisurati e il carico
rimodulato a ogni progresso: gli allenamenti successivi si tarano sugli adattamenti già ottenuti, e
il volume di lavoro resta quantificato, evitando situazioni di sovraccarico.

% ---------------------------------------------------------------------------
\section{Dal test diagnostico alla programmazione}

Sul principio $P = F \times v$, il \textbf{test diagnostico} (\textit{ergo-power}) alimenta cinque
usi distinti:
\begin{itemize}
  \item \textbf{pianificazione e controllo} dei piani di allenamento per atleti di alta prestazione;
  \item \textbf{pianificazione e controllo} dei carichi di lavoro in riabilitazione;
  \item \textbf{pianificazione} dei carichi per l'attività fisica generale;
  \item \textbf{confronto} con i valori standard della popolazione normale;
  \item \textbf{confronto} con i modelli teorico-matematici e ideali della prestazione.
\end{itemize}

\rubrica{La valutazione della capacità di salto}
Gli stessi parametri dell'encoder (spostamento, velocità media, picco di velocità, tempo, forza
media) si applicano al salto, con in più la \textbf{\textit{jump height}}: la differenza tra la
posizione di partenza e quella di arrivo del soggetto. Per il lavoro con i sovraccarichi restano
però decisivi il carico e la percentuale di potenza esprimibile con quel carico.

% ---------------------------------------------------------------------------
\section{Mezzo squat tradizionale e mezzo squat con discesa veloce}

Il \textbf{mezzo squat} si scompone in quattro momenti, letti sulla velocità:
\begin{enumerate}
  \item \textbf{posizione eretta}: velocità nulla;
  \item \textbf{discesa}: la velocità si abbassa rapidamente (valori negativi);
  \item \textbf{massima accosciata}: la velocità torna a zero;
  \item \textbf{inversione} tra fase eccentrica e concentrica e fase finale.
\end{enumerate}
Tra il punto 3 e il punto 4 si colloca il \textbf{picco di velocità} del sollevamento.

\textbf{Mezzo squat con discesa veloce}: variante in cui si velocizza la fase eccentrica, lasciando
invariato il resto dell'esercizio.
\begin{itemize}
  \item accumula più \textbf{energia elastica} e attiva il \textbf{riflesso da stiramento}
    $\rightarrow$ restituiti nella fase concentrica successiva;
  \item modifica i parametri meccanici dell'esercizio e con essi la risposta neuromuscolare.
\end{itemize}

\rubrica{Che cosa cambia nei tracciati}
Registrando forza verticale, posizione verticale del centro di gravità e velocità nelle due
esecuzioni:
\begin{itemize}
  \item \textbf{velocità}: nell'esecuzione veloce la fase eccentrica raggiunge velocità (negative)
    nettamente maggiori e il picco concentrico successivo è più alto; l'intera prova si compie in un
    tempo più breve;
  \item \textbf{potenza}: il picco è più alto nell'esecuzione veloce e viene raggiunto prima ---
    quello che interessa è la \textbf{pendenza della curva}, cioè la capacità di esprimere potenza
    nel minor tempo possibile;
  \item \textbf{forza}: nell'esecuzione veloce il picco è più alto e anticipato rispetto a quella
    tradizionale, sia nella fase di spinta sia all'impatto nella fase di salto.
\end{itemize}
I due sistemi attivano dunque pattern neuromuscolari completamente diversi.

% ---------------------------------------------------------------------------
\section{L'allenamento a isopotenza}

\textbf{Isopotenza}: lavorare a un livello di potenza costante scegliendo carichi diversi. A uno
stesso livello di potenza, sotto il picco della curva potenza-velocità, corrispondono due carichi
distinti --- uno a sinistra e uno a destra.
\begin{itemize}
  \item permette di mantenere il livello di potenza voluto con il carico più basso, senza
    sovraccaricare la schiena;
  \item è il presupposto metodologico del confronto sperimentale che segue.
\end{itemize}

% ---------------------------------------------------------------------------
\section{Il confronto sperimentale tra carico alto e carico basso}

\rubrica{Metodologia}
\begin{itemize}
  \item \textbf{gruppo di controllo}: 5 serie da 6 ripetizioni con il 70\,\% dell'1RM;
  \item \textbf{gruppo sperimentale}: 5 serie da 6 ripetizioni con un carico pari al 45\,\%
    dell'1RM.
\end{itemize}

\rubrica{Risultati}
Dopo 4 settimane di allenamento entrambi i gruppi migliorano, con differenze tra i gruppi per forza,
potenza e sprint ($p < 0{,}05$).

\begin{table}[htbp]
  \centering
  \small
  \renewcommand{\arraystretch}{1.3}
  \begin{tabularx}{\textwidth}{|l|>{\centering\arraybackslash}X|>{\centering\arraybackslash}X|>{\centering\arraybackslash}X|>{\centering\arraybackslash}X|}
    \hline
    & \textbf{RM (kg)} & \textbf{CMJ (cm)} & \textbf{PPO (W)} & \textbf{20 m (s)} \\
    \hline
    \multicolumn{5}{|c|}{\textbf{TG\textsubscript{FAST}} --- carico 45\,\% 1RM} \\
    \hline
    Pre & 120,8 $\pm$ 12,4 & 31,9 $\pm$ 3,6 & 475,1 $\pm$ 66,4 & 3,61 $\pm$ 0,15 \\
    \hline
    Post & 131,9 $\pm$ 18,0 & 36,0 $\pm$ 4,1 & 530,4 $\pm$ 69,9 & 3,54 $\pm$ 0,12 \\
    \hline
    Diff & 11,1 $\pm$ 14 & 4,1 $\pm$ 1,9 & 55,2 $\pm$ 40,5 & $-0,07 \pm 0,10$ \\
    \hline
    \multicolumn{5}{|c|}{\textbf{TG\textsubscript{TRAD}} --- carico 70\,\% 1RM} \\
    \hline
    Pre & 114,2 $\pm$ 14,5 & 30,5 $\pm$ 2,9 & 428,5 $\pm$ 24,8 & 3,64 $\pm$ 0,24 \\
    \hline
    Post & 120,9 $\pm$ 11,6 & 34,5 $\pm$ 3,2 & 459,5 $\pm$ 36,9 & 3,51 $\pm$ 0,15 \\
    \hline
    Diff & 6,7 $\pm$ 16,4 & 4,0 $\pm$ 1,1 & 31,0 $\pm$ 35,6 & $-0,14 \pm 0,12$ \\
    \hline
  \end{tabularx}
\end{table}

Tutte le differenze pre-post sono statisticamente significative in entrambi i gruppi: allenarsi con
un carico molto più basso, purché eseguito in modo esplosivo, produce miglioramenti dello stesso
ordine di quelli ottenuti con carichi alti. La conseguenza metodologica è che si può modulare carico
e volume in funzione delle esigenze del singolo atleta, privilegiando comunque i movimenti esplosivi
e veloci.

% ---------------------------------------------------------------------------
\section{Individualizzazione e specificità del lavoro}

\textbf{Allenamento individuale}: le metodologie devono ruotare attorno alle esigenze del singolo
giocatore, non il contrario --- nel rispetto della sua salute e delle sue potenzialità.
\begin{itemize}
  \item carico e volume si modulano sulle caratteristiche neuromuscolari e metaboliche del singolo
    atleta;
  \item quando il volume di lavoro va ridotto, si privilegiano gli allenamenti che condizionano in
    modo specifico;
  \item il tempo per il lavoro con i sovraccarichi è poco, a fronte dell'enorme carico
    tecnico-tattico e della partita: va reso il più specifico possibile e spendibile il giorno della
    gara.
\end{itemize}

\textbf{Specificità}: coerenza biomeccanica, neuromuscolare e metabolica con i movimenti che si
ritrovano in partita.
\begin{itemize}
  \item in partita nessuno esegue uno squat con i sovraccarichi: ciò che si condiziona non è il
    gesto, ma il \textbf{pattern neuromuscolare} che serve a sprint, accelerazioni, decelerazioni e
    cambi di direzione ad altissima intensità;
  \item il vantaggio è il \textbf{trasferimento} immediato alla prestazione, senza dover attendere
    ulteriore tempo.
\end{itemize}

\rubrica{Il sovraccarico in prossimità della partita}
Con carichi e volumi adeguati il lavoro con i sovraccarichi può essere collocato anche poco prima
della partita:
\begin{itemize}
  \item non come condizionamento, ma come \textbf{riscaldamento e attivazione} dei pattern
    neuromuscolari che serviranno subito dopo;
  \item è sufficiente una serie di 2--3 ripetizioni: più che appesantire il sistema neuromuscolare,
    lo attiva;
  \item con tutte le attenzioni dovute al \textbf{volume}.
\end{itemize}
